%% toc: P2 11 11.8
\section{Proofs for Sheffer stroke, NOR, and expressive limitations}
\label{sec:annex-p02-w02-p02-009-sheffer}

Notation is that of Sections~\ref{sec:p02-w02-p02-008-post}
and~\ref{sec:p02-w02-p02-009-sheffer}.  Besides the arithmetic of
\cref{def:p02-w02-p02-008-bits}, the proofs use three facts about
\(\mathbb B\): \(x\oplus x=0\) and \(xx=x\) for every \(x\); and two members
of \(\mathbb B\) are distinct iff one is \(1\) minus the other.  The
memberships of named operations in the five classes are those of
\cref{thm:p02-w02-p02-008-profile} (Table~\ref{tab:p02-w02-p02-008-profile}).

\begin{annexitem}{annex:p02-w02-p02-009-two-tests}{thm:p02-w02-p02-009-two-tests}{Reversed constants}
\begin{proof}
\emph{1.}  Let \(f(\bar0)=1\) and \(f\in\mathsf M\).  For every
\(\bar a\in\mathbb B^n\), \(\bar0\leq\bar a\), so
\(1=f(\bar0)\leq f(\bar a)\) and \(f(\bar a)=1\); thus \(f=\kappa^n_1\).
Let \(f(\bar1)=0\) and \(f\in\mathsf M\).  For every \(\bar a\),
\(\bar a\leq\bar1\), so \(f(\bar a)\leq f(\bar1)=0\) and \(f(\bar a)=0\);
thus \(f=\kappa^n_0\).

\emph{2.}  Since \(n\geq1\), \(\bar0\leq\bar1\) holds coordinatewise, while
\(f(\bar0)=1>0=f(\bar1)\).  This pair violates the defining condition of
\(\mathsf M\) (\cref{def:p02-w02-p02-008-five}), so \(f\notin\mathsf M\).
(Equivalently, by part~1 a monotone \(f\) with \(f(\bar0)=1\) is
\(\kappa^n_1\) and so has \(f(\bar1)=1\).)

\emph{3.}  Let \(f\in\mathsf L\), with coefficients \(c_0,\ldots,c_n\) such
that \(f(\bar a)=c_0\oplus c_1a_1\oplus\cdots\oplus c_na_n\) for every
\(\bar a\).  At \(\bar a=\bar0\) every product \(c_ia_i\) is \(0\), so
\(c_0=f(\bar0)=1\).  At \(\bar a=\bar1\) the value is
\(1\oplus c_1\oplus\cdots\oplus c_n=f(\bar1)=0\); adding \(1\) to both sides
and using \(1\oplus1=0\) gives \(c_1\oplus\cdots\oplus c_n=1\).  Now let
\(\bar a\) be arbitrary.  The \(i\)-th coordinate of \(\bar a^{\,c}\) is
\(1-a_i=1\oplus a_i\), so by distributivity, commutativity and associativity
\[
 f(\bar a^{\,c})
 =1\oplus\bigoplus_{i=1}^n c_i(1\oplus a_i)
 =1\oplus\Bigl(\bigoplus_{i=1}^n c_i\Bigr)\oplus\bigoplus_{i=1}^n c_ia_i
 =1\oplus1\oplus\bigoplus_{i=1}^n c_ia_i
 =\bigoplus_{i=1}^n c_ia_i .
\]
On the other side, \(f(\bar a)=1\oplus\bigoplus_i c_ia_i\), so
\(1\oplus f(\bar a)=(1\oplus1)\oplus\bigoplus_i c_ia_i=\bigoplus_i c_ia_i\).
Hence \(f(\bar a^{\,c})=1\oplus f(\bar a)=1-f(\bar a)\) for every
\(\bar a\), which is \(f\in\mathsf S\).
\end{proof}
\end{annexitem}

\begin{annexitem}{annex:p02-w02-p02-009-sheffer}{thm:p02-w02-p02-009-sheffer}{Sheffer operations of every arity}
\begin{proof}
\emph{Only if.}  Let \(\langle\{f\}\rangle=\mathcal O\).  Apply
\cref{thm:p02-w02-p02-008-criterion}(1) to \(C=\{f\}\): for each of the five
classes, \(\{f\}\) is not included in it, that is, \(f\) lies outside it.
From \(f\notin\mathsf T_0\), \(f(\bar0)\neq0\), so \(f(\bar0)=1\).  From
\(f\notin\mathsf T_1\), \(f(\bar1)\neq1\), so \(f(\bar1)=0\).  From
\(f\notin\mathsf S\), some \(\bar a\) has \(f(\bar a^{\,c})\neq1-f(\bar a)\);
since \(f(\bar a^{\,c})\) and \(f(\bar a)\) are members of \(\mathbb B\), and
two members are distinct iff one is \(1\) minus the other, this says
\(f(\bar a^{\,c})=f(\bar a)\).

\emph{If.}  Assume conditions 1 to 3.  Condition~1 gives
\(f\notin\mathsf T_0\), condition~2 gives \(f\notin\mathsf T_1\), and
condition~3 gives \(f\notin\mathsf S\), because at the witness \(\bar a\),
\(f(\bar a^{\,c})=f(\bar a)\neq1-f(\bar a)\).  By
\cref{thm:p02-w02-p02-009-two-tests}(2), \(f\notin\mathsf M\).  If
\(f\in\mathsf L\), then \cref{thm:p02-w02-p02-009-two-tests}(3) gives
\(f\in\mathsf S\), contradicting condition~3; so \(f\notin\mathsf L\).  So
for each of the five classes the member \(f\) of \(\{f\}\) lies outside it,
and \cref{thm:p02-w02-p02-008-criterion}(1) gives
\(\langle\{f\}\rangle=\mathcal O\).

\emph{Unary operations.}  For \(n=1\), \(\bar0=(0)\) and \(\bar1=(1)\), so
conditions 1 and 2 force \(f(0)=1\) and \(f(1)=0\), that is,
\(f(a)=1-a\).  Then \(f(1-a)=a=1-f(a)\) for both \(a\), so \(f\) is
self-dual and condition~3 fails.  No unary operation is a Sheffer operation.

\emph{Diagonal.}  The diagonal \(d(x)=f(x,\ldots,x)\) has \(d(0)=f(\bar0)\)
and \(d(1)=f(\bar1)\).  It is negation iff \(d(0)=1\) and \(d(1)=0\), that
is, iff conditions 1 and 2 hold.  For such \(f\) the theorem reduces to the
equivalence of the Sheffer property with condition~3.
\end{proof}
\end{annexitem}

\begin{annexitem}{annex:p02-w02-p02-009-count}{thm:p02-w02-p02-009-count}{Sheffer count}
\begin{proof}
Fix \(n\geq1\).  By \cref{thm:p02-w02-p02-009-sheffer} the Sheffer
operations in \(\mathcal O_n\) are the members of
\(A=\{f\in\mathcal O_n:f(\bar0)=1,\ f(\bar1)=0\}\) that are not in
\(\mathsf S\).  So their number is \(|A|-|A\cap\mathsf S|\).

\emph{The set \(A\).}  An operation in \(\mathcal O_n\) is an arbitrary
assignment of a value in \(\mathbb B\) to each of the \(2^n\) rows.  As
\(n\geq1\), the rows \(\bar0\) and \(\bar1\) are distinct.  Members of \(A\)
have fixed values on these two rows and arbitrary values on the other
\(2^n-2\) rows, so \(|A|=2^{2^n-2}\).

\emph{Self-dual operations.}  Let \(R\) be the set of the \(2^{n-1}\) rows
\(\bar a\) with \(a_1=0\).  For every row \(\bar a\), exactly one of
\(\bar a\) and \(\bar a^{\,c}\) lies in \(R\), because their first
coordinates are \(a_1\) and \(1-a_1\).  Restriction to \(R\) is a bijection
from \(\mathsf S\cap\mathcal O_n\) onto the set of all functions
\(R\to\mathbb B\).  It is injective: if \(f,f'\in\mathsf S\) agree on \(R\)
and \(\bar a\notin R\), then \(\bar a^{\,c}\in R\), and
\(f(\bar a)=f((\bar a^{\,c})^{c})=1-f(\bar a^{\,c})=1-f'(\bar a^{\,c})
=f'(\bar a)\).  It is surjective: given \(u:R\to\mathbb B\), put
\(f(\bar a)=u(\bar a)\) for \(\bar a\in R\) and
\(f(\bar a)=1-u(\bar a^{\,c})\) for \(\bar a\notin R\).  For
\(\bar a\in R\), \(\bar a^{\,c}\notin R\) and
\(f(\bar a^{\,c})=1-u(\bar a)=1-f(\bar a)\); for \(\bar a\notin R\),
\(\bar a^{\,c}\in R\) and \(f(\bar a^{\,c})=u(\bar a^{\,c})=1-f(\bar a)\).
So \(f\in\mathsf S\), and its restriction to \(R\) is \(u\).  Hence
\(|\mathsf S\cap\mathcal O_n|=2^{2^{n-1}}\).

\emph{The set \(A\cap\mathsf S\).}  For \(f\in\mathsf S\),
\(f(\bar1)=f(\bar0^{\,c})=1-f(\bar0)\), so \(f\in A\) iff \(f(\bar0)=1\).
The row \(\bar0\) lies in \(R\).  Under the bijection above, the members of
\(A\cap\mathsf S\) correspond to the functions \(u:R\to\mathbb B\) with
\(u(\bar0)=1\), of which there are \(2^{2^{n-1}-1}\).

\emph{Count and values.}  The number of Sheffer operations is
\(2^{2^n-2}-2^{2^{n-1}-1}\).  For \(n=1\) it is \(2^0-2^0=0\); for \(n=2\),
\(2^2-2^1=2\); for \(n=3\), \(2^6-2^3=64-8=56\); for \(n=4\),
\(2^{14}-2^7=16384-128=16256\).

\emph{Proportion.}  Dividing by \(|\mathcal O_n|=2^{2^n}\) and using
\(2^n=2^{n-1}+2^{n-1}\),
\[
 \frac{2^{2^n-2}-2^{2^{n-1}-1}}{2^{2^n}}
 =2^{-2}-2^{2^{n-1}-1-2^n}
 =\tfrac14-2^{-2^{n-1}-1}.
\]
As \(n\) increases, the exponent \(-2^{n-1}-1\) strictly decreases and tends
to \(-\infty\), so the subtracted term strictly decreases and tends to
\(0\); the proportion strictly increases and tends to \(\tfrac14\).
\end{proof}
\end{annexitem}

\begin{annexitem}{annex:p02-w02-p02-009-single}{thm:p02-w02-p02-009-single}{Single connectives}
\begin{proof}
Throughout, positive-arity completeness of a connective set \(C\) is decided
by \cref{thm:p02-w02-p02-008-criterion}(2): it holds iff \(C^+\) has, for
each of the five classes, a member outside it.  By
Table~\ref{tab:p02-w02-p02-008-profile}, \(\kappa^1_0\) lies in
\(\mathsf T_0\), \(\mathsf M\) and \(\mathsf L\) and outside \(\mathsf T_1\)
and \(\mathsf S\); \(\kappa^1_1\) lies in \(\mathsf T_1\), \(\mathsf M\) and
\(\mathsf L\) and outside \(\mathsf T_0\) and \(\mathsf S\).

\emph{1.}  Let \(c\) be nullary with value \(b\).  By
\cref{def:p02-w02-p02-008-arity}, \(\{c\}^+=\{\kappa^1_b\}\), and
\(\kappa^1_b\in\mathsf M\).  So no member of \(\{c\}^+\) lies outside
\(\mathsf M\), and \(\{c\}\) is not positive-arity functionally complete.
By \cref{thm:p02-w02-p02-008-nullary}(2), nullary-inclusive completeness
requires positive-arity completeness, so \(\{c\}\) is not nullary-inclusive
functionally complete either.

\emph{2.}  Here \(\{c\}^+=\{f\}\).  By
\cref{thm:p02-w02-p02-008-criterion}, \(\{c\}\) is positive-arity
functionally complete iff \(\langle\{f\}\rangle=\mathcal O\), that is, iff
\(f\) is a Sheffer operation, and by \cref{thm:p02-w02-p02-009-sheffer} iff
\(f\) satisfies the three conditions there.  The set \(\{c\}\) has no
nullary member, so by \cref{thm:p02-w02-p02-008-nullary}(1) it is not
nullary-inclusive functionally complete.

\emph{3.}  A connective set with exactly one member is \(\{c\}\) with \(c\)
nullary or of positive arity.  By parts~1 and~2 it is not nullary-inclusive
functionally complete.  By \cref{def:p02-w02-p02-008-completeness},
functional completeness in the sense of
Section~\ref{sec:prop-syntax-semantics-04} is nullary-inclusive
completeness.

\emph{4.}  Each of the three sets contains \(\bottom\) or \(\topv\), a
nullary member, so by \cref{thm:p02-w02-p02-008-nullary}(2) it is
nullary-inclusive functionally complete iff it is positive-arity
functionally complete.  It remains to decide positive-arity completeness.

\emph{(a)}  \(\{c,\bottom\}^+=\{f,\kappa^1_0\}\).  The member
\(\kappa^1_0\) lies outside \(\mathsf T_1\) and \(\mathsf S\) and inside
\(\mathsf T_0\), \(\mathsf M\) and \(\mathsf L\).  So the set is complete iff
\(f\notin\mathsf T_0\), \(f\notin\mathsf M\) and \(f\notin\mathsf L\).  If
the set is complete, then \(f\notin\mathsf T_0\), so \(f(\bar0)=1\), and
\(f\notin\mathsf L\).  Conversely, let \(f(\bar0)=1\) and
\(f\notin\mathsf L\).  Then \(f\notin\mathsf T_0\).  If \(f\in\mathsf M\),
then \cref{thm:p02-w02-p02-009-two-tests}(1) gives \(f=\kappa^n_1\), which
lies in \(\mathsf L\) (take \(c_0=1\) and \(c_1=\cdots=c_n=0\) in
\cref{def:p02-w02-p02-008-five}); this contradicts \(f\notin\mathsf L\).  So
\(f\notin\mathsf M\), and the set is complete.

\emph{(b)}  \(\{c,\topv\}^+=\{f,\kappa^1_1\}\).  The member \(\kappa^1_1\)
lies outside \(\mathsf T_0\) and \(\mathsf S\) and inside \(\mathsf T_1\),
\(\mathsf M\) and \(\mathsf L\).  So the set is complete iff
\(f\notin\mathsf T_1\), \(f\notin\mathsf M\) and \(f\notin\mathsf L\).
Completeness gives \(f(\bar1)=0\) and \(f\notin\mathsf L\).  Conversely, if
\(f(\bar1)=0\) and \(f\notin\mathsf L\), then \(f\notin\mathsf T_1\), and
\(f\in\mathsf M\) would give \(f=\kappa^n_0\) by
\cref{thm:p02-w02-p02-009-two-tests}(1), which is affine with all
coefficients \(0\); so \(f\notin\mathsf M\) and the set is complete.

\emph{(c)}  \(\{c,\bottom,\topv\}^+=\{f,\kappa^1_0,\kappa^1_1\}\).  Here
\(\kappa^1_0\) lies outside \(\mathsf T_1\) and \(\mathsf S\), and
\(\kappa^1_1\) outside \(\mathsf T_0\); both lie in \(\mathsf M\) and
\(\mathsf L\).  So the set is complete iff \(f\notin\mathsf M\) and
\(f\notin\mathsf L\).

\emph{Closed formulas.}  The formula \(c(\bottom,\ldots,\bottom)\), with
\(n\) arguments, contains no variable, and by the truth clauses its value
under every valuation is \(f(0,\ldots,0)=f(\bar0)\), which is \(1\) in
case~(a); so it represents the nullary constant \(1\), and \(\bottom\)
represents \(0\).  Likewise \(c(\topv,\ldots,\topv)\) has value
\(f(\bar1)=0\) in case~(b), and \(\topv\) represents \(1\).

\emph{5.}  If \(c\) is a Sheffer connective, then
\(\langle\{f\}\rangle=\mathcal O\), and by
\cref{thm:p02-w02-p02-008-criterion}(1) \(f\) lies outside each of the five
classes.  So \(f(\bar0)=1\), \(f(\bar1)=0\), \(f\notin\mathsf M\) and
\(f\notin\mathsf L\), and the conditions of (a), (b) and (c) all hold.
\end{proof}
\end{annexitem}

\begin{annexitem}{annex:p02-w02-p02-009-ternary}{thm:p02-w02-p02-009-ternary}{Ternary instances}
\begin{proof}
Write \(\bar a=(x,y,z)\).  A clone \(K\) with \(C\subseteq K\) satisfies
\(\langle C\rangle\subseteq K\) by \cref{thm:p02-w02-p02-008-generated}(5),
and \(\mathsf S\) and \(\mathsf L\) are clones by
\cref{thm:p02-w02-p02-008-five-clones}.

\emph{1.}  Let \(N(x,y,z)=1\oplus xyz\).  Then
\(N(x,y,y)=1\oplus xyy=1\oplus xy=\mathrm{NAND}(x,y)\), using \(yy=y\); so
\(\mathrm{NAND}=N\circ(e^2_1,e^2_2,e^2_2)\in\langle\{N\}\rangle\).  By
\cref{thm:p02-w02-p02-008-nand-nor}(2) and
\cref{thm:p02-w02-p02-008-generated}(3),
\(\mathcal O=\langle\{\mathrm{NAND}\}\rangle\subseteq\langle\{N\}\rangle\).
Let \(R(x,y,z)=(1\oplus x)(1\oplus y)(1\oplus z)\).  Then
\(R(x,y,y)=(1\oplus x)(1\oplus y)(1\oplus y)=(1\oplus x)(1\oplus y)
=\mathrm{NOR}(x,y)\), and in the same way
\(\mathcal O=\langle\{\mathrm{NOR}\}\rangle\subseteq\langle\{R\}\rangle\).

\emph{2.}  \(\operatorname{maj}(\bar0)=0\) and \(\operatorname{maj}(\bar1)=1\),
so \(g(\bar0)=1\) and \(g(\bar1)=0\); hence \(g\notin\mathsf T_0\),
\(g\notin\mathsf T_1\), and \(g\notin\mathsf M\) by
\cref{thm:p02-w02-p02-009-two-tests}(2).
\emph{Self-duality.}  If \(\bar a\) has \(k\) coordinates equal to \(1\),
then \(\bar a^{\,c}\) has \(3-k\).  Now \(k\geq2\) iff \(3-k\leq1\) iff not
\(3-k\geq2\); so \(\operatorname{maj}(\bar a^{\,c})=1-\operatorname{maj}(\bar a)\),
and \(g(\bar a^{\,c})=1-\operatorname{maj}(\bar a^{\,c})
=\operatorname{maj}(\bar a)=1-g(\bar a)\).  So \(g\in\mathsf S\).
\emph{Not affine.}  By \cref{thm:p02-w02-p02-008-tests}(3), if
\(g\in\mathsf L\) its coefficients are \(c_0=g(\bar0)=1\) and
\(c_i=g(\bar0[i{:=}1])\oplus1\).  Each \(\bar0[i{:=}1]\) has one coordinate
\(1\), so \(\operatorname{maj}\) is \(0\) there, \(g\) is \(1\), and
\(c_i=0\).  The affine form would be the constant \(1\), but
\(g(\bar1)=0\); so \(g\notin\mathsf L\).
\emph{Consequences.}  \(\{g\}\subseteq\mathsf S\) gives
\(\langle\{g\}\rangle\subseteq\mathsf S\).  The operations \(\land\),
\(\lor\), \(\to\), \(\mathrm{NAND}\) and \(\mathrm{NOR}\) lie outside
\(\mathsf S\) (Table~\ref{tab:p02-w02-p02-008-profile}), and so does every
constant operation \(\kappa^n_b\), since
\(\kappa^n_b(\bar a^{\,c})=b\neq1-b=1-\kappa^n_b(\bar a)\).  So none of
them lies in \(\langle\{g\}\rangle\).  If a connective set \(C\) has
\(C^+=\{g\}\), every formula over \(C\) represents a member of
\(\langle C^+\rangle\subseteq\mathsf S\) by
\cref{thm:p02-w02-p02-008-formulas-clone}, so no formula over \(C\)
represents any of these operations.

\emph{3.}  \(h(\bar a)=1\oplus x\oplus y\oplus z\) has the form of
\cref{def:p02-w02-p02-008-five} with \(c_0=c_1=c_2=c_3=1\), so
\(h\in\mathsf L\).  \(h(\bar0)=1\), and \(h(\bar1)=1\oplus1\oplus1\oplus1=0\)
because \(1\oplus1=0\) twice.  So \(h\notin\mathsf T_0\),
\(h\notin\mathsf T_1\), \(h\notin\mathsf M\) by
\cref{thm:p02-w02-p02-009-two-tests}(2), and \(h\in\mathsf S\) by
\cref{thm:p02-w02-p02-009-two-tests}(3).  Since \(\{h\}\subseteq\mathsf S\)
and \(\{h\}\subseteq\mathsf L\), \(\langle\{h\}\rangle\) is included in both,
so in \(\mathsf S\cap\mathsf L\).  The operations excluded in part~2 lie
outside \(\mathsf S\), and \(\operatorname{maj}\) lies outside \(\mathsf L\)
(Table~\ref{tab:p02-w02-p02-008-profile}); none of them lies in
\(\langle\{h\}\rangle\), and as in part~2 no formula over a connective set
with operation set \(\{h\}\) represents any of them.

\emph{4.}  Recall \((u\to v)=1\oplus u\oplus uv\), which is \(0\) iff
\(u=1\) and \(v=0\).  Then \(f(x,x,x)=(1-x)\cdot(x\to x)=(1-x)\cdot1=1-x\).
Further \(f(0,1,0)=1\cdot(1\to0)=0\) and \(f(1,0,1)=0\cdot(0\to1)=0\), so
\(f(x,1-x,x)=0\) for both values of \(x\); and
\(f(0,y,z)=1\cdot(y\to z)=(y\to z)\).  For the Sheffer property:
\(f(\bar0)=1\cdot(0\to0)=1\); \(f(\bar1)=0\cdot(1\to1)=0\); and at
\(\bar a=(0,1,0)\), \(f(\bar a)=0=f(1,0,1)=f(\bar a^{\,c})\), so \(f\) is
not self-dual.  By \cref{thm:p02-w02-p02-009-sheffer}, \(f\) is a Sheffer
operation.

\emph{The definitions of the worked example.}  These are the formulas of
\cref{def:p02-w02-p02-009-worked}.  By the truth clauses, the formula
\(c(p_1,p_1,p_1)\) has the truth function \(x\mapsto f(x,x,x)=1-x\).  The
formula
\(\theta=c(p_1,c(p_1,p_1,p_1),p_1)\) has truth function
\(x\mapsto f(x,1-x,x)=0\).  The formula \(\iota=c(\theta,p_1,p_2)\) has
truth function \((x,y)\mapsto f(0,x,y)=(x\to y)\), because the value of
\(\theta\) is \(0\) at every row.  The formula \(\iota'\), obtained by
putting \(c(p_2,p_2,p_2)\) for \(p_2\) in \(\iota\), has value
\(x\to(1-y)\) at \((x,y)\), by the substitution fact (F2) of
Section~\ref{sec:annex-p02-w02-p02-008-post}.  Finally
\(c(\iota',\iota',\iota')\) has value \(1-(x\to(1-y))\), which is \(1\) iff
\(x\to(1-y)\) is \(0\), iff \(x=1\) and \(1-y=0\), iff \(x=y=1\); so it has
truth function \(\land\).
\end{proof}
\end{annexitem}
